\section{Applications of the theorem on formal functions}
\label{section-applications-formal-functions}


\noindent
We will add more here as needed. For the moment we need the
following characterization of finite morphisms in the Noetherian case.

\begin{lemma}
\label{lemma-characterize-finite}
(For a more general version see
More on Morphisms, Lemma \ref{more-morphisms-lemma-characterize-finite}.)
Let $f : X \to S$ be a morphism of schemes.
Assume $S$ is locally Noetherian.
The following are equivalent
\begin{enumerate}
\item $f$ is finite, and
\item $f$ is proper with finite fibres.
\end{enumerate}
\end{lemma}

\begin{proof}
A finite morphism is proper according to
Morphisms, Lemma \ref{morphisms-lemma-finite-proper}.
A finite morphism is quasi-finite according to
Morphisms, Lemma \ref{morphisms-lemma-finite-quasi-finite}.
A quasi-finite morphism has finite fibres, see
Morphisms, Lemma \ref{morphisms-lemma-quasi-finite}.
Hence a finite morphism is proper and has finite fibres.

\medskip\noindent
Assume $f$ is proper with finite fibres.
We want to show $f$ is finite.
In fact it suffices to prove $f$ is affine.
Namely, if $f$ is affine, then it follows that
$f$ is integral by
Morphisms, Lemma \ref{morphisms-lemma-integral-universally-closed}
whereupon it follows from
Morphisms, Lemma \ref{morphisms-lemma-finite-integral}
that $f$ is finite.

\medskip\noindent
To show that $f$ is affine we may assume that $S$ is affine, and our
goal is to show that $X$ is affine too.
Since $f$ is proper we see that $X$ is separated and quasi-compact.
Hence we may use the criterion of
Lemma \ref{lemma-quasi-separated-h1-zero-covering} to prove that $X$
is affine. To see this let $\mathcal{I} \subset \mathcal{O}_X$
be a finite type ideal sheaf. In particular $\mathcal{I}$ is
a coherent sheaf on $X$. By
Lemma \ref{lemma-higher-direct-images-zero-finite-fibre} we conclude that
$R^1f_*\mathcal{I}_s = 0$ for all $s \in S$.
In other words, $R^1f_*\mathcal{I} = 0$. Hence we see from
the Leray Spectral Sequence for $f$ that
$H^1(X , \mathcal{I}) = H^1(S, f_*\mathcal{I})$.
Since $S$ is affine, and $f_*\mathcal{I}$ is quasi-coherent
(Schemes, Lemma \ref{schemes-lemma-push-forward-quasi-coherent})
we conclude $H^1(S, f_*\mathcal{I}) = 0$
from Lemma \ref{lemma-quasi-coherent-affine-cohomology-zero}
as desired. Hence $H^1(X, \mathcal{I}) = 0$ as desired.
\end{proof}

\noindent
As a consequence we have the following useful result.

\begin{lemma}
\label{lemma-proper-finite-fibre-finite-in-neighbourhood}
\begin{slogan}
A proper morphism is finite in a neighbourhood of a finite fiber.
\end{slogan}
(For a more general version see
More on Morphisms,
Lemma \ref{more-morphisms-lemma-proper-finite-fibre-finite-in-neighbourhood}.)
Let $f : X \to S$ be a morphism of schemes.
Let $s \in S$.
Assume
\begin{enumerate}
\item $S$ is locally Noetherian,
\item $f$ is proper, and
\item $f^{-1}(\{s\})$ is a finite set.
\end{enumerate}
Then there exists an open neighbourhood $V \subset S$ of $s$
such that $f|_{f^{-1}(V)} : f^{-1}(V) \to V$ is finite.
\end{lemma}

\begin{proof}
The morphism $f$ is quasi-finite at all the points of $f^{-1}(\{s\})$
by Morphisms, Lemma \ref{morphisms-lemma-finite-fibre}.
By Morphisms, Lemma \ref{morphisms-lemma-quasi-finite-points-open} the
set of points at which $f$ is quasi-finite is an open $U \subset X$.
Let $Z = X \setminus U$. Then $s \not \in f(Z)$. Since $f$ is proper
the set $f(Z) \subset S$ is closed. Choose any open neighbourhood
$V \subset S$ of $s$ with $Z \cap V = \emptyset$. Then
$f^{-1}(V) \to V$ is locally quasi-finite and proper.
Hence it is quasi-finite
(Morphisms, Lemma \ref{morphisms-lemma-quasi-finite-locally-quasi-compact}),
hence has finite fibres
(Morphisms, Lemma \ref{morphisms-lemma-quasi-finite}), hence
is finite by Lemma \ref{lemma-characterize-finite}.
\end{proof}

\begin{lemma}
\label{lemma-ample-on-fibre}
Let $f : X \to Y$ be a proper morphism of schemes with $Y$ Noetherian.
Let $\mathcal{L}$ be an invertible $\mathcal{O}_X$-module.
Let $\mathcal{F}$ be a coherent $\mathcal{O}_X$-module.
Let $y \in Y$ be a point such that $\mathcal{L}_y$ is ample on $X_y$.
Then there exists a $d_0$ such that for all $d \geq d_0$ we have
$$
R^pf_*(\mathcal{F} \otimes_{\mathcal{O}_X} \mathcal{L}^{\otimes d})_y = 0
\text{ for }p > 0
$$
and the map
$$
f_*(\mathcal{F} \otimes_{\mathcal{O}_X} \mathcal{L}^{\otimes d})_y
\longrightarrow
H^0(X_y, \mathcal{F}_y \otimes_{\mathcal{O}_{X_y}} \mathcal{L}_y^{\otimes d})
$$
is surjective.
\end{lemma}

\begin{proof}
Note that $\mathcal{O}_{Y, y}$ is a Noetherian local ring.
Consider the canonical morphism
$c : \Spec(\mathcal{O}_{Y, y}) \to Y$, see
Schemes, Equation (\ref{schemes-equation-canonical-morphism}).
This is a flat morphism as it identifies local rings.
Denote momentarily $f' : X' \to \Spec(\mathcal{O}_{Y, y})$
the base change of $f$ to this local ring. We see that
$c^*R^pf_*\mathcal{F} = R^pf'_*\mathcal{F}'$ by
Lemma \ref{lemma-flat-base-change-cohomology}.
Moreover, the fibres $X_y$ and $X'_y$ are identified.
Hence we may assume that $Y = \Spec(A)$ is the spectrum of
a Noetherian local ring $(A, \mathfrak m, \kappa)$ and $y \in Y$
corresponds to $\mathfrak m$. In this case
$R^pf_*(\mathcal{F} \otimes_{\mathcal{O}_X} \mathcal{L}^{\otimes d})_y =
H^p(X, \mathcal{F} \otimes_{\mathcal{O}_X} \mathcal{L}^{\otimes d})$
for all $p \geq 0$. Denote $f_y : X_y \to \Spec(\kappa)$ the projection.

\medskip\noindent
Let $B = \text{Gr}_\mathfrak m(A) =
\bigoplus_{n \geq 0} \mathfrak m^n/\mathfrak m^{n + 1}$.
Consider the sheaf $\mathcal{B} = f_y^*\widetilde{B}$
of quasi-coherent graded $\mathcal{O}_{X_y}$-algebras.
We will use notation as in Section \ref{section-theorem-formal-functions}
with $I$ replaced by $\mathfrak m$.
Since $X_y$ is the closed subscheme of $X$ cut out by
$\mathfrak m\mathcal{O}_X$ we may think of
$\mathfrak m^n\mathcal{F}/\mathfrak m^{n + 1}\mathcal{F}$
as a coherent $\mathcal{O}_{X_y}$-module, see
Lemma \ref{lemma-i-star-equivalence}. Then
$\bigoplus_{n \geq 0} \mathfrak m^n\mathcal{F}/\mathfrak m^{n + 1}\mathcal{F}$
is a quasi-coherent graded $\mathcal{B}$-module of finite type
because it is generated in degree zero over $\mathcal{B}$
and because the degree zero part is
$\mathcal{F}_y = \mathcal{F}/\mathfrak m \mathcal{F}$
which is a coherent $\mathcal{O}_{X_y}$-module.
Hence by Lemma \ref{lemma-graded-finiteness} part (2)
we see that
$$
H^p(X_y, \mathfrak m^n \mathcal{F}/ \mathfrak m^{n + 1}\mathcal{F}
\otimes_{\mathcal{O}_{X_y}} \mathcal{L}_y^{\otimes d}) = 0
$$
for all $p > 0$, $d \geq d_0$, and $n \geq 0$. By
Lemma \ref{lemma-relative-affine-cohomology}
this is the same as the statement that
$
H^p(X, \mathfrak m^n \mathcal{F}/ \mathfrak m^{n + 1}\mathcal{F}
\otimes_{\mathcal{O}_X} \mathcal{L}^{\otimes d}) = 0
$
for all $p > 0$, $d \geq d_0$, and $n \geq 0$.

\medskip\noindent
Consider the short exact sequences
$$
0 \to \mathfrak m^n\mathcal{F}/\mathfrak m^{n + 1} \mathcal{F}
\to \mathcal{F}/\mathfrak m^{n + 1} \mathcal{F}
\to \mathcal{F}/\mathfrak m^n \mathcal{F} \to 0
$$
of coherent $\mathcal{O}_X$-modules. Tensoring with $\mathcal{L}^{\otimes d}$
is an exact functor and we obtain short exact sequences
$$
0 \to
\mathfrak m^n\mathcal{F}/\mathfrak m^{n + 1} \mathcal{F}
\otimes_{\mathcal{O}_X} \mathcal{L}^{\otimes d}
\to \mathcal{F}/\mathfrak m^{n + 1} \mathcal{F}
\otimes_{\mathcal{O}_X} \mathcal{L}^{\otimes d}
\to \mathcal{F}/\mathfrak m^n \mathcal{F}
\otimes_{\mathcal{O}_X} \mathcal{L}^{\otimes d} \to 0
$$
Using the long exact cohomology sequence and the vanishing above
we conclude (using induction) that
\begin{enumerate}
\item $H^p(X, \mathcal{F}/\mathfrak m^n \mathcal{F}
\otimes_{\mathcal{O}_X} \mathcal{L}^{\otimes d}) = 0$
for all $p > 0$, $d \geq d_0$, and $n \geq 0$, and
\item $H^0(X, \mathcal{F}/\mathfrak m^n \mathcal{F}
\otimes_{\mathcal{O}_X} \mathcal{L}^{\otimes d}) \to
H^0(X_y, \mathcal{F}_y \otimes_{\mathcal{O}_{X_y}} \mathcal{L}_y^{\otimes d})$
is surjective for all $d \geq d_0$ and $n \geq 1$.
\end{enumerate}
By the theorem on formal functions (Theorem \ref{theorem-formal-functions})
we find that the $\mathfrak m$-adic completion of
$H^p(X, \mathcal{F} \otimes_{\mathcal{O}_X} \mathcal{L}^{\otimes d})$
is zero for all $d \geq d_0$ and $p > 0$.
Since $H^p(X, \mathcal{F} \otimes_{\mathcal{O}_X} \mathcal{L}^{\otimes d})$
is a finite $A$-module by
Lemma \ref{lemma-proper-over-affine-cohomology-finite}
it follows from Nakayama's lemma (Algebra, Lemma \ref{algebra-lemma-NAK})
that $H^p(X, \mathcal{F} \otimes_{\mathcal{O}_X} \mathcal{L}^{\otimes d})$
is zero for all $d \geq d_0$ and $p > 0$.
For $p = 0$ we deduce from
Lemma \ref{lemma-ML-cohomology-powers-ideal} part (3)
that $H^0(X, \mathcal{F} \otimes_{\mathcal{O}_X} \mathcal{L}^{\otimes d}) \to
H^0(X_y, \mathcal{F}_y \otimes_{\mathcal{O}_{X_y}} \mathcal{L}_y^{\otimes d})$
is surjective, which gives the final statement of the lemma.
\end{proof}

\begin{lemma}
\label{lemma-ample-in-neighbourhood}
(For a more general version see
More on Morphisms,
Lemma \ref{more-morphisms-lemma-ample-in-neighbourhood}.)
Let $f : X \to Y$ be a proper morphism of schemes with $Y$ Noetherian.
Let $\mathcal{L}$ be an invertible $\mathcal{O}_X$-module.
Let $y \in Y$ be a point such that $\mathcal{L}_y$ is ample
on $X_y$. Then there is an open neighbourhood $V \subset Y$
of $y$ such that $\mathcal{L}|_{f^{-1}(V)}$ is ample on $f^{-1}(V)/V$.
\end{lemma}

\begin{proof}
Pick $d_0$ as in Lemma \ref{lemma-ample-on-fibre} for
$\mathcal{F} = \mathcal{O}_X$. Pick $d \geq d_0$
so that we can find $r \geq 0$ and sections
$s_{y, 0}, \ldots, s_{y, r} \in H^0(X_y, \mathcal{L}_y^{\otimes d})$
which define a closed immersion
$$
\varphi_y =
\varphi_{\mathcal{L}_y^{\otimes d}, (s_{y, 0}, \ldots, s_{y, r})} :
X_y \to \mathbf{P}^r_{\kappa(y)}.
$$
This is possible by Morphisms, Lemma
\ref{morphisms-lemma-finite-type-over-affine-ample-very-ample}
but we also use
Morphisms, Lemma \ref{morphisms-lemma-image-proper-scheme-closed}
to see that $\varphi_y$ is a closed immersion and
Constructions, Section \ref{constructions-section-projective-space}
for the description of morphisms into projective
space in terms of invertible sheaves and sections.
By our choice of $d_0$, after replacing $Y$ by an open neighbourhood
of $y$, we can choose
$s_0, \ldots, s_r \in H^0(X, \mathcal{L}^{\otimes d})$
mapping to $s_{y, 0}, \ldots, s_{y, r}$.
Let $X_{s_i} \subset X$ be the open subset where $s_i$
is a generator of $\mathcal{L}^{\otimes d}$. Since
the $s_{y, i}$ generate $\mathcal{L}_y^{\otimes d}$ we see that
$X_y \subset U = \bigcup X_{s_i}$.
Since $X \to Y$ is closed, we see that
there is an open neighbourhood $y \in V \subset Y$
such that $f^{-1}(V) \subset U$.
After replacing $Y$ by $V$ we may assume that
the $s_i$ generate $\mathcal{L}^{\otimes d}$. Thus we
obtain a morphism
$$
\varphi = \varphi_{\mathcal{L}^{\otimes d}, (s_0, \ldots, s_r)} :
X \longrightarrow \mathbf{P}^r_Y
$$
with $\mathcal{L}^{\otimes d} \cong \varphi^*\mathcal{O}_{\mathbf{P}^r_Y}(1)$
whose base change to $y$ gives $\varphi_y$.

\medskip\noindent
We will finish the proof by a sleight of hand; the ``correct'' proof
proceeds by directly showing that $\varphi$ is a closed
immersion after base changing to an open neighbourhood of $y$.
Namely, by Lemma \ref{lemma-proper-finite-fibre-finite-in-neighbourhood}
we see that $\varphi$ is a finite over an open neighbourhood
of the fibre $\mathbf{P}^r_{\kappa(y)}$ of $\mathbf{P}^r_Y \to Y$
above $y$. Using that $\mathbf{P}^r_Y \to Y$ is closed, after
shrinking $Y$ we may assume that $\varphi$ is finite.
Then $\mathcal{L}^{\otimes d} \cong \varphi^*\mathcal{O}_{\mathbf{P}^r_Y}(1)$
is ample by the very general
Morphisms, Lemma \ref{morphisms-lemma-pullback-ample-tensor-relatively-ample}.
\end{proof}





